\section*{अध्याय \ref{ch:g10:universal-gravitation} --- सार्वत्रिक गुरुत्वाकर्षण और भार}

\begin{solution}{exo:g10:universal-gravitation:1}
\emph{1.} $F = \num{6.67e-11} \times \dfrac{60 \times 60}{0.80^2}
= \dfrac{\num{2.40e-7}}{0.64} \approx \qty{3.8e-7}{N}$।

\emph{2.} मच्छर का भार $P = \num{2.5e-6} \times 9.81 \approx \qty{2.5e-5}{N}$ है --- यानी नर्तकों के पारस्परिक
आकर्षण से लगभग $65$ गुना \emph{बड़ा}। लोगों के बीच का गुरुत्व एक मच्छर
के भार से भी कमज़ोर है।
\end{solution}

\begin{solution}{exo:g10:universal-gravitation:2}
\emph{1.} $P = mg = 55 \times 9.81 \approx \qty{5.4e2}{N}$।

\emph{2.} उसका द्रव्यमान अपरिवर्तित रहता है: \qty{55}{kg} (पदार्थ की वही
मात्रा)। उसका भार $P = 55 \times 1.62 \approx \qty{89}{N}$ हो जाता है --- यानी छह गुना कम।
\end{solution}

\begin{solution}{exo:g10:universal-gravitation:3}
\emph{1.} $g_{\text{मंगल}} = \dfrac{G M}{R^2}
= \dfrac{\num{6.67e-11} \times \num{6.42e23}}{(\num{3.39e6})^2}
= \dfrac{\num{4.28e13}}{\num{1.15e13}} \approx \qty{3.73}{N/kg}$।

\emph{2.} मंगल पर: $P = 900 \times 3.73 \approx \qty{3.4e3}{N}$। पृथ्वी पर: $P = 900 \times 9.81 \approx \qty{8.8e3}{N}$ --- द्रव्यमान वही रहते हुए
रोवर मंगल पर लगभग $2.6$ गुना कम भारी है।
\end{solution}

\begin{solution}{exo:g10:universal-gravitation:4}
\[
F = \num{6.67e-11} \times
\frac{\num{5.97e24} \times \num{7.35e22}}{(\num{3.84e8})^2}
= \frac{\num{2.93e37}}{\num{1.47e17}}
\approx \qty{2.0e20}{N}.
\]
चंद्रमा पृथ्वी पर \emph{उतने ही} परिमाण का, यानी \qty{2.0e20}{N} का बल विपरीत
दिशा में लगाता है: \omterm{def:g10:universal-gravitation:force}{गुरुत्वाकर्षण} परस्पर होता है।
\end{solution}

\begin{solution}{exo:g10:universal-gravitation:5}
\emph{1.} दूरी दुगुनी: $d^2$ $4$ से गुणा हो जाता है, इसलिए $F = F_0/4$।

\emph{2.} दूरी आधी: $F = 4F_0$।

\emph{3.} दोनों द्रव्यमान दुगुने: $m_1 m_2$ $4$ से गुणा हो जाता है,
इसलिए $F = 4F_0$।

\emph{4.} अंश $\times 3$, हर $\times 9$: $F = F_0/3$।
\end{solution}

\begin{solution}{exo:g10:universal-gravitation:6}
\emph{1.} $F = \num{6.67e-11} \times
\dfrac{(\num{5.0e8})^2}{100^2}
= \num{6.67e-11} \times \dfrac{\num{2.5e17}}{\num{1.0e4}}
\approx \qty{1.7e3}{N}$।

\emph{2.} $m = F/g = 1667/9.81 \approx \qty{1.7e2}{kg}$: दोनों जहाज़ एक-दूसरे को एक बड़ी मोटरसाइकिल के
भार जितने बल से खींचते हैं।

\emph{3.} $a = F/m = \num{1.7e3}/\num{5.0e8} \approx
\qty{3.3e-6}{m/s^2}$। कोई विरोध न हो तब भी जहाज़ों को पैदल चलने की चाल पकड़ने
में लगभग पाँच दिन लगेंगे; व्यवहार में पानी का प्रतिरोध और लंगर इस आकर्षण
को पूरी तरह दबा देते हैं।
\end{solution}

\begin{solution}{exo:g10:universal-gravitation:7}
\emph{1.} दंड तुला दोनों लोकों पर \qty{3.0}{kg} बताती है: थैली और मानक
द्रव्यमान, दोनों के भार उसी गुणक से भाग हो जाते हैं, इसलिए तुलना
जस की तस रहती है। कमानीदार तराज़ू भार पढ़ता है: पृथ्वी पर $3.0 \times 9.81 \approx \qty{29}{N}$
और चंद्रमा पर $3.0 \times 1.62 \approx \qty{4.9}{N}$।

\emph{2.} तुला: वह द्रव्यमान नापती है, और ग्राहक वही ख़रीद रहा है।
पृथ्वी पर ``\omterm{def:g10:orders-of-magnitude:unit}{किलोग्राम}'' में अंशांकित कमानीदार तराज़ू उसी थैली के
लिए चंद्रमा पर $4.9/9.81 \approx \qty{0.50}{kg}$ पढ़ेगा --- और दुकानदार एक की क़ीमत में छह थैलियाँ
दे बैठेगा।
\end{solution}

\begin{solution}{exo:g10:universal-gravitation:8}
\emph{1.} केंद्र से दूरी: $d = R + h = \num{6.37e6} + \num{4.0e5} = \qty{6.77e6}{m}$। तब $g = \dfrac{\num{6.67e-11} \times \num{5.97e24}}{(\num{6.77e6})^2}
= \dfrac{\num{3.98e14}}{\num{4.58e13}} \approx \qty{8.69}{N/kg}$।

\emph{2.} सतह वाले मान का $8.69/9.81 \approx 89\%$।

\emph{3.} वहाँ ऊपर गुरुत्व लगभग पूरी ताक़त पर है: अंतरिक्ष यात्री इसलिए
तैरते हैं कि स्टेशन और उसमें रखा सब कुछ पृथ्वी के चारों ओर \emph{साथ-साथ
गिर} रहे हैं (मुक्त गिरावट), इसलिए नहीं कि गुरुत्व मिट गया है।
\end{solution}

\begin{solution}{exo:g10:universal-gravitation:9}
$g = GM/R^2$ से,
\[
M = \frac{g R^2}{G}
= \frac{8.87 \times (\num{6.05e6})^2}{\num{6.67e-11}}
= \frac{\num{3.25e14}}{\num{6.67e-11}}
\approx \qty{4.87e24}{kg},
\]
यानी पृथ्वी के द्रव्यमान का लगभग $0.8$ गुना --- शुक्र आकार और
द्रव्यमान में लगभग पृथ्वी का जुड़वाँ है।
\end{solution}

\begin{solution}{exo:g10:universal-gravitation:10}
\emph{1.} सूर्य: $F = \dfrac{\num{6.67e-11} \times \num{1.99e30}}{(\num{1.496e11})^2}
\approx \qty{5.9e-3}{N}$। चंद्रमा: $F = \dfrac{\num{6.67e-11} \times \num{7.35e22}}{(\num{3.84e8})^2}
\approx \qty{3.3e-5}{N}$।

\emph{2.} सूर्य लगभग $\num{5.9e-3}/\num{3.3e-5} \approx 180$ गुना अधिक ज़ोर से खींचता है।

\emph{3.} पृथ्वी का व्यास पृथ्वी--चंद्रमा दूरी का कहीं बड़ा \emph{अंश}
($\approx 3\%$) है, जबकि पृथ्वी--सूर्य दूरी का बहुत छोटा ($\approx 0.01\%$), इसलिए
चंद्रमा का खिंचाव गोले के पास वाले पहलू से दूर वाले पहलू तक कहीं अधिक
बदलता है। ज्वार इसी अंतर पर पलते हैं --- और वहाँ चंद्रमा जीत जाता है।
\end{solution}

\begin{solution}{exo:g10:universal-gravitation:11}
\emph{1.} $g' = \dfrac{G (2M)}{(2R)^2} = \dfrac{2}{4}\,\dfrac{GM}{R^2}
= \dfrac{g}{2} \approx \qty{4.9}{N/kg}$: दुगुनी त्रिज्या दुगुने द्रव्यमान पर भारी पड़ती है।

\emph{2.} हमें $\dfrac{GM}{R'^2} = 2\,\dfrac{GM}{R^2}$ चाहिए, इसलिए $R'^2 = R^2/2$ और $R' = R/\sqrt{2} \approx \qty{4.5e6}{m}$ --- यानी लगभग
\qty{4500}{km}।
\end{solution}

\begin{solution}{exo:g10:universal-gravitation:12}
\emph{1.} प्रति \omterm{def:g10:orders-of-magnitude:unit}{किलोग्राम} खिंचाव बराबर: $\dfrac{G M_{\text{पृथ्वी}}}{x^2} = \dfrac{G M_{\text{चंद्रमा}}}{(D-x)^2}$। तिरछा गुणा करके
और $G$ से भाग देकर: $\left(\dfrac{x}{D-x}\right)^2 = \dfrac{M_{\text{पृथ्वी}}}{M_{\text{चंद्रमा}}}$।

\emph{2.} $\dfrac{x}{D-x}
= \sqrt{\dfrac{\num{5.97e24}}{\num{7.35e22}}} = \sqrt{81.2} \approx
9.0$, इसलिए $x = 9.0\,(D - x)$, यानी $x = \dfrac{9.0}{10.0}\,D
\approx 0.90 \times \num{3.84e8} \approx \qty{3.5e8}{m}$। संतुलन-बिंदु चंद्रमा तक की
यात्रा के $90\%$ पर बैठता है: लगभग पूरी यात्रा पर पृथ्वी का ही खिंचाव
हावी रहता है।
\end{solution}

\begin{solution}{exo:g10:universal-gravitation:13}
\emph{1.} फैलाने पर: $R^2 + 2Rh + h^2 = R^2 + x^2$; चूँकि $h$ $R$ की तुलना में नन्हा है,
$h^2$ छोड़ दीजिए: $2Rh \approx x^2$, इसलिए $h \approx \dfrac{x^2}{2R}
= \dfrac{(\num{8.0e3})^2}{2 \times \num{6.37e6}}
= \dfrac{\num{6.4e7}}{\num{1.27e7}} \approx \qty{5.0}{m}$।

\emph{2.} एक \omterm{def:g10:orders-of-magnitude:unit}{सेकंड} में प्रक्षेप्य \qty{4.9}{m} गिरता है --- लगभग ठीक उतना
\qty{5.0}{m}, जितना ज़मीन \qty{8.0}{km} पर नीचे हट जाती है। इसलिए $v \approx \qty{8.0}{km/s}$ पर गिरावट
ज़मीन को कभी पकड़ नहीं पाती: प्रक्षेप्य कक्षा में आ जाता है।

\emph{3.} $T = \dfrac{2\pi R}{v}
= \dfrac{2\pi \times \num{6.37e6}}{\num{8.0e3}}
\approx \qty{5.0e3}{s} \approx \qty{83}{min}$ --- जो असली निचली कक्षा वाले उपग्रहों
के लगभग $\qty{90}{min}$ के पास है, और वे कुछ ऊँचे उड़ते हैं।
\end{solution}

\begin{solution}{exo:g10:universal-gravitation:14}
\emph{1.} $g = \dfrac{\num{6.67e-11} \times \num{2.8e30}}
{(\num{1.2e4})^2} = \dfrac{\num{1.87e20}}{\num{1.44e8}}
\approx \qty{1.3e12}{N/kg}.$

\emph{2.} $P = 70 \times \num{1.3e12} \approx \qty{9.1e13}{N}$, जबकि पृथ्वी पर $\qty{6.9e2}{N}$: यानी लगभग $\num{1.3e11}$ गुना अधिक।
परमाणुओं से बनी कोई संरचना वहाँ खड़ी नहीं रह सकती।

\emph{3.} $P = \num{1.0e-6} \times \num{1.3e12} \approx
\qty{1.3e6}{N}$। पृथ्वी पर यह $m = \num{1.3e6}/9.81 \approx \qty{1.3e5}{kg}$ का भार है --- यानी रेत का
एक कण, और भार सौ टन के रेल इंजन जितना।
\end{solution}

\begin{solution}{exo:g10:universal-gravitation:15}
\emph{1.} दूरी $60$ गुना बड़ी है, इसलिए प्रति \omterm{def:g10:orders-of-magnitude:unit}{किलोग्राम}
खिंचाव $60^2 = 3600$ गुना कमज़ोर है।

\emph{2.} एक \omterm{def:g10:orders-of-magnitude:unit}{सेकंड} की गिरावट भी उसी तरह बदलती है: $4.9/3600 \approx \qty{1.4e-3}{m}$ --- यानी लगभग
\qty{1.4}{mm}।

\emph{3.} $T = 27.3 \times \qty{86400}{s} = \qty{2.36e6}{s}$, इसलिए $v = \dfrac{2\pi \times \num{3.84e8}}{\num{2.36e6}} \approx
\qty{1.02e3}{m/s}$, और एक \omterm{def:g10:orders-of-magnitude:unit}{सेकंड} में $x = \qty{1.02e3}{m}$। गिरावट
$h \approx \dfrac{x^2}{2d}
= \dfrac{(\num{1.02e3})^2}{2 \times \num{3.84e8}}
\approx \qty{1.4e-3}{m}$ है।

\emph{4.} दोनों संख्याएँ मेल खाती हैं: चंद्रमा हर \omterm{def:g10:orders-of-magnitude:unit}{सेकंड} पृथ्वी की ओर ठीक
उतना ही गिरता है जितना गिरते सेब पर अंशांकित व्युत्क्रम-वर्ग गुरुत्व
बताता है। न्यूटन ने निष्कर्ष निकाला कि जो बल चंद्रमा को थामे है वही बल
सेब को गिराता है --- गुरुत्व सार्वत्रिक है।
\end{solution}

\begin{solution}{pb:g10:universal-gravitation:1}
\textbf{1.} $F = \num{6.67e-11} \times
\dfrac{\num{7.3e6} \times 55}{100^2} \approx \qty{2.7e-6}{N}$।

\textbf{2.} रेत का कण $\num{1.0e-6} \times 9.81 \approx \qty{9.8e-6}{N}$ भारी है: पूरी एफ़िल मीनार पर्यटक को रेत के
एक कण के भार से लगभग $3.7$ गुना \emph{कम} खींचती है।

\textbf{3.} (क) दूरी $\times 10$: बल $100$ से भाग। (ख) दोनों द्रव्यमान
$\times 1000$: बल $10^6$ से गुणा। (ग) द्रव्यमान $\times 1000$ और दूरी $\times 1000$:
$10^6 / 10^6$ --- बल अपरिवर्तित।

\textbf{4.} $F = \dfrac{\num{6.67e-11} \times \num{5.97e24} \times
1.00}{(\num{6.37e6})^2} \approx \qty{9.81}{N}$: एक \omterm{def:g10:orders-of-magnitude:unit}{किलोग्राम} का भार --- और यही
ठीक-ठीक क्षेत्र प्रबलता $g = \qty{9.81}{N/kg}$ है।

\textbf{5.} गुरुत्व प्रति \omterm{def:g10:orders-of-magnitude:unit}{किलोग्राम} कमज़ोर है, पर वह हमेशा
केवल \emph{आकर्षित} करता है, इसलिए पृथ्वी के सारे $\num{5.97e24}$
\omterm{def:g10:orders-of-magnitude:unit}{किलोग्राम} के खिंचाव एक ही दिशा में जुड़ जाते हैं --- और खगोलीय
द्रव्यमान पड़ोसी वाले द्रव्यमान को दस की बाईस घातों से हरा देता है।

\textbf{6.} $F = \num{6.67e-11} \times \dfrac{158 \times 0.73}{0.23^2}
= \num{6.67e-11} \times \dfrac{115.3}{0.0529}
\approx \qty{1.5e-7}{N}$।

\textbf{7.} छोटा गोला $0.73 \times 9.81 \approx \qty{7.2}{N}$ भारी है --- यानी जिस आकर्षण को पकड़ना है
उससे लगभग $\num{5e7}$ गुना अधिक। पलड़े वाला कोई तराज़ू पाँच करोड़ में एक
हिस्से का विभेदन नहीं कर सकता; पर बारीक मरोड़ तार, जो नैनोन्यूटन स्तर
के बल-आघूर्णों पर भी दिखने भर को ऐंठ जाता है, कर सकता है।

\textbf{8.} $G = \dfrac{F d^2}{m_1 m_2}
= \dfrac{\num{1.47e-7} \times 0.0529}{115.3}
\approx \qty{6.74e-11}{N.m^2/kg^2}$ --- यानी 1798 में एक लकड़ी के छप्पर से निकला मान,
जो आधुनिक \num{6.67e-11} के लगभग $1\%$ के भीतर है।

\textbf{9.} $g = G M / R^2$ से:
\[
M = \frac{g R^2}{G}
= \frac{9.81 \times (\num{6.37e6})^2}{\num{6.67e-11}}
\approx \qty{5.97e24}{kg}.
\]
साठ लाख अरब अरब \omterm{def:g10:orders-of-magnitude:unit}{किलोग्राम}: पृथ्वी, तुल गई।

\textbf{10.} $V = \frac43 \pi R^3 = \frac43 \pi
(\num{6.37e6})^3 \approx \qty{1.08e21}{m^3}$, इसलिए $\rho = \num{5.97e24}/\num{1.08e21} \approx \qty{5.5e3}{kg/m^3}$ --- यानी सतही चट्टान के घनत्व का
दुगुना। भीतर ग्रेनाइट से कहीं सघन कुछ छिपा होना चाहिए: पृथ्वी का एक
भारी (लोहे का) \omterm{def:g10:refraction:fiber}{क्रोड} है।

\textbf{11.} $g_{\text{चंद्रमा}} = \dfrac{\num{6.67e-11} \times
\num{7.35e22}}{(\num{1.74e6})^2} = \dfrac{\num{4.90e12}}
{\num{3.03e12}} \approx \qty{1.62}{N/kg}$।

\textbf{12.} ऊँचाई $\propto 1/g$: $h = 0.50 \times \dfrac{9.81}{1.62} \approx \qty{3.0}{m}$ --- यानी खड़े-खड़े बास्केटबॉल के
घेरे के ऊपर से छलाँग, वह भी धीमी गति में।

\textbf{13.} $g_{\text{बुध}} = \dfrac{\num{6.67e-11} \times
\num{3.30e23}}{(\num{2.44e6})^2} = \dfrac{\num{2.20e13}}
{\num{5.95e12}} \approx \qty{3.70}{N/kg}$ --- लगभग ठीक मंगल का \qty{3.73}{N/kg}। बुध का द्रव्यमान कम
है, पर उसकी छोटी त्रिज्या (नीचे छोटा $R^2$) इसकी भरपाई कर देती है: वह
कहीं अधिक सघन लोक है।

\textbf{14.} आधी क्षेत्र प्रबलता के लिए $(R + h)^2 = 2 R^2$ चाहिए, इसलिए $R + h = R\sqrt{2}$ और
$h = (\sqrt{2} - 1) R \approx 0.414 \times \num{6.37e6}
\approx \qty{2.6e6}{m}$ --- यानी लगभग \qty{2600}{km} ऊपर।

\textbf{15.} $\dfrac{G M_{\text{पृथ्वी}}}{d^2} = 1.62$ से $d = \sqrt{\dfrac{\num{6.67e-11} \times \num{5.97e24}}{1.62}}
= \sqrt{\num{2.46e14}} \approx \qty{1.57e7}{m} \approx 2.5$ पृथ्वी-त्रिज्याएँ मिलती हैं। ज़मीन से
केवल डेढ़ त्रिज्या ऊपर से ही पृथ्वी चंद्रमा की सतह से अधिक ज़ोर से नहीं
खींचती।

\textbf{16.} $v = \dfrac{2\pi d}{T}
= \dfrac{2\pi \times \num{1.496e11}}{\num{3.156e7}}
\approx \qty{2.98e4}{m/s}$ --- यानी हर \omterm{def:g10:orders-of-magnitude:unit}{सेकंड} तीस किलोमीटर।

\textbf{17.} $\dfrac{v^2}{d} = \dfrac{(\num{2.98e4})^2}{\num{1.496e11}}
\approx \qty{5.9e-3}{N/kg}$: सूर्य पृथ्वी के हर \omterm{def:g10:orders-of-magnitude:unit}{किलोग्राम} को लगभग छह
मिलीन्यूटन से खींचता है।

\textbf{18.} $\dfrac{G M_{\text{सूर्य}}}{d^2} = \dfrac{v^2}{d}$ से
\[
M_{\text{सूर्य}} = \frac{v^2 d}{G}
= \frac{\num{8.87e8} \times \num{1.496e11}}{\num{6.67e-11}}
\approx \qty{1.99e30}{kg}.
\]

\textbf{19.} $\dfrac{\num{1.99e30}}{\num{5.97e24}} \approx
\num{3.3e5}$: सूर्य तीन लाख पृथ्वियों के बराबर है।

\textbf{20.} चूँकि वही $G$ द्रव्यमानों के हर युग्म पर राज करता है,
इसलिए उसे एक छप्पर में सीसे के दो गोलों के बीच एक बार नाप लेने से हम
पहले अपने पैरों तले के ग्रह को और फिर उस तारे को तराज़ू पर रख सके
जिसकी हम कक्षा में हैं --- \emph{सार्वत्रिक} शब्द की क़ीमत
यही है।
\end{solution}
